\documentclass[aspectratio=43]{beamer}
\usepackage{amsmath,amsfonts}
\usepackage{algorithmic}
\usepackage{algorithm}
\usepackage{array}
\usepackage{caption}
%\usepackage[caption=false,font=normalsize,labelfont=sf,textfont=sf]{subfig}
\usepackage{mathtools}
\usepackage{pgfpages}
\usepackage{textcomp}
\usepackage{stfloats}
\usepackage{url}
\usepackage{verbatim}
\usepackage{graphicx}
\usepackage{cite}
\usetheme{metropolis}
\captionsetup{font=tiny,labelformat=empty,justification=centering}
\setlength{\abovecaptionskip}{1pt}

\title{Micro Reentry Detection :)}
\date{November 25, 2024}
\author{Phil Alcorn, Sadie Shudde}
\institute{Tarleton State University}

\begin{document}
%\pgfpagesuselayout{4 on 1}[letterpaper,border shrink=5mm]

\maketitle

\begin{frame}{Table of Contents}
\tableofcontents
\end{frame}

\section{Overview}

\begin{frame}{Overview of Micro Reentry}
\begin{itemize}
    \item \textbf{Definition:} Micro reentry circuits are small-scale reentrant 
		pathways in excitable media (e.g., cardiac tissue) that sustain 
		self-propagating electrical activity.\cite{aha_microreentry2016}

    \item \textbf{Mechanism:}
    \begin{itemize}
        \item Requires a closed conduction loop
        \item Unidirectional block in part of the pathway
        \item Conduction delay allowing tissue ahead to recover excitability
    \end{itemize}

	\item \textbf{Characteristic Scale:} 
    \begin{itemize}
        \item Typically on the order of millimeters or smaller
        \item Often below spatial resolution of clinical imaging
    \end{itemize}

    \item \textbf{Physiological Context:}
    \begin{itemize}
        \item Commonly observed in cardiac arrhythmias (e.g., atrial 
			fibrillation, ventricular tachycardia)
        \item Associated with heterogeneous tissue properties (fibrosis, ischemia)
    \end{itemize}

\end{itemize}
\vspace{-5pt}

\end{frame}

\begin{frame}{Basic Reentry}
\begin{itemize}
    \item A reentry circuit is a circulating excitation wave in a closed pathway of excitable tissue.

    \item \textbf{Structure:} Loop (anatomical or functional) with cell-to-cell 
		conduction. \cite{biorxiv_microreentry2021}

    \item \textbf{Sustaining Condition:}
    \begin{itemize}
        \item Wave returns after tissue has recovered excitability
        \item Requires sufficient path length relative to conduction delay and refractory period
    \end{itemize}

    \item \textbf{Key Dynamics:}
    \begin{itemize}
        \item Finite conduction velocity
        \item Refractory period limits re-excitation
		\item More likely to occur in more heterogeneus tissue
		\item Often represented as a 1D ring of excitable elements with a circulating wave.
    \end{itemize}

\end{itemize}
\vspace{-8pt}
\end{frame}


\begin{frame}{Micro Reentry vs Macro Reentry}
\begin{itemize}
	\item \textbf{Macro-reentry:} \cite{cvphysiology_reentry} 
    \begin{itemize}
		\item Large-scale circuit around anatomical structures (e.g., valves, scars, \textbf{pulmonary veins})
        \item Path length is long relative to wavelength ($\lambda = v \cdot \text{ERP}$)
        \item Typically mappable with clinical techniques (e.g., ECG, catheter mapping)
    \end{itemize}

	\item \textbf{Micro-reentry:} \cite{oliveira2018microreentry}
    \begin{itemize}
        \item Small, localized circuits driven by tissue heterogeneity
        \item Path length comparable to wavelength (near minimum for sustained reentry)
        \item Often below clinical resolution; inferred via modeling or high-density mapping
    \end{itemize}
\end{itemize}
\end{frame}

\begin{frame}{Why Study Reentry?}

\end{frame}

\begin{frame}{Papers discussing micro reentry after ablation}


\end{frame}

\begin{frame}{Papers talking about the size of micro circuits}


\end{frame}

\section{What we are doing}

\begin{frame}{Our Goal}

\end{frame}

\begin{frame}{How does this goal fit with the other projects we have?}
    
\end{frame}

\begin{frame}{Simple mesh program}


\end{frame}

\section{Results}
\begin{frame}{Interesting values so far }
	Open questions of why would muscles have significantly smaller RPs 
\end{frame}

\begin{frame}{Transferring setups from the mini-sim to full model}
    
\end{frame}
\begin{frame}{Current Limitations}
    
\end{frame}
\begin{frame}{Future plans for the project}
   Hand it over to Dr. Wyatt 
\end{frame}
\begin{frame}[allowframebreaks]{References}
	\bibliographystyle{IEEEtran}
	\bibliography{references}
\end{frame}
\end{document}
